% !TEX program = xelatex
\documentclass[12pt,oneside]{ctexart}

% ---- 页面布局与边注 ----
\usepackage[
  a4paper,
  left=50mm,           % 左边留白，留出边注区
  right=25mm,
  top=25mm,
  bottom=28mm,
  marginparwidth=36mm, % 边注栏宽度
  marginparsep=6mm     % 边注与正文间距
]{geometry}
\setlength{\parskip}{0.4em}
\setlength{\parindent}{2em}

% 边注：固定在左侧
\usepackage{marginnote}
\renewcommand*{\marginfont}{\small\itshape}
\reversemarginpar
\newcommand{\n}[2][0em]{\marginnote{\raggedright #2}[#1]}

% ---- 数学与结构 ----
\usepackage{amsmath,amssymb,amsthm}
\usepackage{enumitem}
\setlist{nosep}

\theoremstyle{definition}
\newtheorem{definition}{定义}[section]
\theoremstyle{plain}
\newtheorem{theorem}[definition]{定理}
\theoremstyle{remark}
\newtheorem{example}[definition]{例}

% ---- 页眉页脚 ----
\usepackage{fancyhdr}
\pagestyle{fancy}
\fancyhf{}
\lhead{\footnotesize 天体物理中的辐射过程学习笔记}
\rhead{\footnotesize \leftmark}
\cfoot{\thepage}

% ---- 颜色与超链接 ----
\usepackage{xcolor}
\definecolor{link}{RGB}{17,85,204}
\usepackage[
  colorlinks=true,
  linkcolor=link,
  urlcolor=link,
  citecolor=link
]{hyperref}

% ---- 附加推导环境（可自定义标题 & 可一键隐藏）----
\newif\ifshowderiv
\showderivtrue                 % 全局开关：隐藏时改为 \showderivfalse

\newcommand{\derivdefaulttitle}{额外说明}

\newenvironment{derivationnote}[1][\derivdefaulttitle]{%
  % 开始：开启跨环境条件，false 时整段内容都会被跳过
  \ifshowderiv\begin{quote}\small\itshape
  \noindent\textbf{#1}\;\par
}{%
  \end{quote}\fi
}


% ---- 章节标题 ----
\usepackage{titlesec}
\titleformat{\section}{\Large\bfseries}{\thesection}{0.8em}{}
\titleformat{\subsection}{\large\bfseries}{\thesubsection}{0.6em}{}

% ---- 文档信息 ----
\title{天体物理中的辐射过程学习笔记}
\author{Leonhard Hsiao}
\date{\today}

\begin{document}
\maketitle
\tableofcontents
\section{同步辐射和回旋辐射}


回旋辐射和同步辐射本质上都是带电粒子在磁场中受洛伦兹力作用做加速运动时产生的电磁辐射。根据电子的速度不同，辐射的特性也有所不同。
\begin{itemize}
  \item 非相对论电子产生的辐射叫回旋辐射（Cyclotron radiation）
  \item 相对论电子产生的辐射叫同步辐射（Synchrotron radiation）
\end{itemize}

\textbf{比较轫致辐射：} 无磁场等离子中的自由电子在离子库伦场作用下获得加速度而产生的辐射。

\begin{definition}
  \textbf{回旋辐射：} 非相对论电子在磁场中受洛伦兹力作用做加速运动时产生的辐射。其辐射频率低，功率小。
\end{definition}

\n{只有强磁场中的天体，回旋辐射才不可忽视，在太阳耀斑、白矮星的光学辐射、中子星的X射线辐射等强磁场天体中，回旋辐射不可忽视。
}
\subsection{相对论性电子在磁场中的运动}
四维洛伦兹力为：
\begin{align}
  F^\mu = \frac{q}{c}F^{\mu\nu}U_\nu = \frac{1}{c}\gamma
  \begin{pmatrix}
    \vec{E}\cdot\vec{v}/c \\
    \vec{E}c + \vec{v}\times\vec{B}
  \end{pmatrix}
\end{align}
由此可列的带电粒子在电磁场中的运动方程：
\begin{align}
  \frac{dP^{\mu}}{d\tau} = \frac{d}{d\tau} \gamma m 
  \begin{pmatrix}
    c \\
    \vec{v}
  \end{pmatrix} = 
  \frac{q}{c}
  \begin{pmatrix}
    \vec{E} \cdot \vec{v} \\
    \vec{E}c + \vec{v} \times \vec{B}
  \end{pmatrix}
\end{align}
诚然，带电粒子的运动会产生辐射损失能量，辐射阻尼会扭曲螺旋装的电子运动轨迹，但是同步辐射能量极高，在考虑全局辐射的情况下，
同步辐射在极短的时标内释放能量，所以可以忽略辐射阻尼对电子轨迹的影响。所以磁场中粒子运动方程可以忽略辐射能量变换。
从另外一个角度说一下辐射阻尼，在电子辐射电磁场会反向影响整体的电磁场分布，电子自己的运动会受到自己产生磁场的阻碍，这就是辐射阻尼。
非相对论性粒子产生的电磁场不足以影响原电磁场，而高能粒子产生的电磁场虽然可以影响电磁场，但是同步辐射产生的时标极短，之后电子就退化为非相对论情况，所以分析同步辐射也可以忽略辐射阻尼。

    
\subsection{同步辐射}
同步辐射本质上还是带电粒子做加速运动产生的辐射，可以用相对论性电磁的辐射功率求得同步辐射，
\begin{align}
  P = \frac{2q^2}{3c^3}\gamma^4(a_\perp^2 + \gamma^2 a_\parallel^2)
\end{align}
在磁场中的电子加速度垂直于速度，所以有：
\begin{align}
  p = \frac{2q^2}{3c^3}\gamma^2\frac{q^2B^2}{m^2c^2}v^2\sin^2\alpha
\end{align}
其中，$\alpha$ 是电子速度与磁场方向的夹角，称为\textbf{跃迁角}。对于电子，其特殊情况为：
\begin{align}
  P = \frac{2}{3}r_0^{2}cB^2\gamma^2\beta^2\sin^2\alpha = 1.6\times 10^{-15}B^2\gamma^2\beta^2\sin^2\alpha \text{ erg s}^{-1}
\end{align}
\n{注意粒子能量是 $E = \gamma mc^2$，磁场能量密度是 $U_B = B^2/8\pi$}
其中，$r_0 = e^2/mc^2 = 2.82\times 10^{-13}$ cm 是经典电子半径。所以，电子的同步辐射功率与洛伦兹因子成正比，与磁场能量密度成正比。
如果取非相对论极限，即 $\gamma \to 1$，$\beta \leqq 1$则上式退化为回旋辐射功率公式。
\begin{align}
  p = \frac{2}{3}\frac{e^4}{m^2c^5}v^2B^2\sin^2\alpha
\end{align}
不难发现，由于洛伦兹因子的存在，同步辐射的功率以$\gamma^2$倍大于回旋辐射功率。单电子辐射在各个方向上的平均是：
\begin{align}
  \langle P \rangle = \frac{4}{3}\sigma_T c U_B \gamma^2 \beta^2\gamma^2U_B
\end{align}
其中，$\sigma_T = (8\pi/3)r_0^2$ 是汤姆逊散射截面。
\end{document}